\documentclass[a4paper]{article}
    \usepackage{fullpage}
    \usepackage{amsmath}
    \usepackage{amssymb}
    \usepackage{textcomp}
    \usepackage[utf8]{inputenc}
    \usepackage{hyperref}
    \usepackage[T1]{fontenc}
    \textheight=10in
    \pagestyle{empty}
    \raggedright
    \usepackage[left=0.8in,right=0.8in,bottom=0.8in,top=0.8in]{geometry}

\def\bull{\vrule height 0.8ex width .7ex depth -.1ex }

% DEFINITIONS FOR RESUME %%%%%%%%%%%%%%%%%%%%%%%

\newcommand{\lineunder} {
    \vspace*{-8pt} \\
    \hspace*{-18pt} \hrulefill \\
}

\newcommand{\header} [1] {
    {\hspace*{-18pt}\vspace*{6pt} \textsc{#1}}
    \vspace*{-6pt} \lineunder
}

\newcommand{\contact} [3] {
    \vspace*{20pt}
    \begin{center}
        {\Huge \textbf{#1}}\\
        \vspace{5pt}
        #2 \\ #3
    \end{center}
    \vspace*{-8pt}
}

\usepackage{xcolor}
\definecolor{cvblue}{RGB}{0, 50, 140} % Professional Navy

\usepackage{hyperref}
\hypersetup{
    colorlinks=true,
    linkcolor=black,     % Keep internal page links black
    urlcolor=cvblue,      % Make external evidence links blue
    pdfborder={0 0 0}
}

% END RESUME DEFINITIONS %%%%%%%%%%%%%%%%%%%%%%%

\begin{document}
\vspace*{-40pt}

%==== Profile ====%
\vspace*{-10pt}
\vspace*{6mm}
\begin{center}
    {\Huge \textbf{Aaromal A.}}\\
    \vspace{8pt}
    Kochi, India $\cdot$ aaromalonline@gmail.com $\cdot$ \href{https://www.linkedin.com/in/aaromalonline/}{https://www.linkedin.com/in/aaromalonline/} \\
    \vspace{6pt}
    Electronics Engineering $\cdot$ AI/ML Research 
\end{center}
\vspace*{-8pt}

%==== Education ====%
\vspace*{9mm}
\header{Education}
\textbf{\href{https://www.cbse.gov.in/}{Central Board of Secondary Education (CBSE)}} \hfill Kerala, 2022 -- 2024\\
\textit{Higher Secondary Education, Mathematics and Computer Science} \hfill Grade: 94.6\%\\
\textbf{Coursework:} Python Programming, Database Management Systems (DBMS), Data Structures and Algorithms (DSA), Computer Networking, Statistics \& Probability, Coordinate Geometry, Calculus.\\
\textbf{Activities:} Kerala State Junior Meet Athlete, District Junior Athletic Meet Medalist, School Cadre Group Leader.
\vspace{4mm}

\textbf{\href{https://www.cusat.ac.in/}{Cochin University of Science \& Technology (CUSAT)}} \hfill Kerala, 2024 -- 2028\\
\textit{B.Tech in Electronics \& Communication Engineering} \hfill CGPA: 8.67 (S3)\\
\textbf{Major Coursework:} Analog Integrated Circuits, Digital System Design, Signals and Systems, Network Theory, Microprocessors \& Microcontrollers (8086, 8051), C programming \& Object-Oriented Programming (C++), Linear Algebra \& Transform Techniques, Numerical \& Statistical Methods, Calculus.
\vspace*{4mm}

%==== Research & Presentations ====%
\header{Research \& Presentations}
\textbf{\href{https://github.com/aaromalonline/ecometachip}{Ecometachip Mini: Low-Cost, Eco-friendly, Dielectric Chemical Identifier}} \hfill \textit{New Delhi} \\
\textit{\href{https://epapers2.org/apscon2026/ESR/paper_details.php?paper_id=6101}{Lecture Presentation, IEEE APSCON 2026 (Student Research Forum)}} \hfill \textit{23-25 Feb 2026}
\begin{itemize} \itemsep 1pt
    \item Designed a eco-friendly dielectric sensing system for preliminary identification of common liquids (water, ethanol, acetone, petrol) based on relative permittivity variation in a biodegradable coir-rubber composite as the primary sensing element.
    \item Developed a 100 MHz Hartley oscillator-based readout circuit with RF pickup for magnetic field modulation proportional to variation in rel. permittivity of material with different liquids and precision rectifier for rectification in following stages.
    \item Implemented LED-based RGB visual classification to validate output voltage variations (low cost).
    \item Simulated ideal behavior in \textbf{LTspice} by varying coupling constants and analyzed the sensitivity and repeatability of the sensor under variable environmental conditions (Humidity/Temp) using \textbf{Python (Pandas/Seaborn)}.
    \item Quantified environmental drift coefficients ($mV/\%RH$ and $mV/^\circ C$) to characterize sensor robustness and output consistency across variable ambient conditions.
    \item Prepared research paper using \LaTeX\ which got selected for lecture presentation under \href{https://epapers2.org/apscon2026/ESR/session_view.php?session_id=35}{Student Research Forum}, \textbf{\href{https://2026.ieee-apscon.org/}{IEEE APSCON'26} (IEEE Applied Sensing Conference)} by \href{https://ieee-sensors.org/}{IEEE Sensors Council} held at New Delhi, 23-25 Feb, 2026. (\href{https://www.researchgate.net/publication/401928480_Certificate_for_presentation}{ResearchGate presentation DOI})
\end{itemize}

\textbf{Attention-Regularized Vision Transformer for Cloth-Agnostic Person Re-Identification.} \hfill \textit{Kottayam} \\
\textit{Ongoing Research, \href{https://internship.iiitkottayam.ac.in/}{Institute Internship Programme}, IIIT Kottayam} \hfill \textit{May 2026 -- Present}
\begin{itemize} \itemsep 1pt
    \item Researching cloth-agnostic person re-identification (\textbf{CC-ReID, PU-ReID}) under \textbf{\href{https://scholar.google.com/citations?user=dGwfp0wAAAAJ&hl=en}{Dr.\ Sridhar Raj S}}, begun as a summer research internship under the IIITK Institute Internship Programme.
    \item Writing, training, and evaluating \textbf{Vision Transformer (ViT-Base, PyTorch)} models with regularized attention to improve robustness to clothing changes.
    \item Conducted multiple ablation studies, achieving improvements over current benchmark models on the task. Writing a journal paper out of it and in process on publishing on a journal.
\end{itemize}
\vspace*{2mm}

%==== Featured Projects ====%
{\hspace*{-18pt}\vspace*{6pt} \textsc{Featured Projects} \hfill \href{https://github.com/aaromalonline}{github.com/aaromalonline}}
\vspace*{-6pt} \lineunder

\textbf{\href{https://liberate-atcs.vercel.app/}{Liberate:ATCS - Adaptive Typing \& Control System}} \hfill \textit{Feb 2025 - Sep 2025}\\
{\sl ADXL345, TCRT5000 IR, ESP32, C++, Python, Arduino IDE}
\begin{itemize} \itemsep 1pt
    \item Implimented muscle-twitch decoding interface using \textbf{ADXL345 accelerometer} (also \textbf{TCTRT5000 IR sensor}) as input modules to record muscle twitches as 3-axis accelerations, which are encoded into binary clicks using baseline calibration and filtering algorithms.
    \item Data is sent serially/wirelessly via ESP32 to the \textbf{Python application} using I2C, where the user can control a keyboard interface using the decoded twitch clicks.
    \item The \textbf{PyQt5 interface} provides a moving highlight bar which spans across the rows of a adaptive keyboard interface (incl QWERTY, SOS, Mobility, Speak Keys), allowing selection of rows and keys through muscle twitches.
    \item  The decoded twitch signals can be used for various applications incl communication, mobility (used to move wheelchair in addition to a motor-contolller) etc, aiming to assist people with ALS or paralysis.
    \item Received participation recognition at the \textbf{\href{https://sscs.ieee.org/education/sscs-arduino-contest/}{2025 IEEE SSCS Arduino Content}} and \textbf{\href{https://stridedesignathon.ieee-link.org/}{STRIDE Assitive Designation 2025}} (Social Technology \& Research for Inclusive Design Excellence) by \href{https://kdisc.kerala.gov.in/en/}{K-DISC} (Kerala Development and Innovation Strategic Council) documenting \textbf{\href{https://youtu.be/9NmNVjKmtew?si=vOsgkYYBrDFDhQ74}{video presentation}} and demonstration of the prototype.
\end{itemize}

\href{https://github.com/aaromalonline/cardiokey}{\textbf{CardioKey (ECGID) - Analog ECG Acquisition \& Biometrics}} \hfill \textit{Feb 2026 - Present}\\
{\sl AD8232 ECG \& ADS1115 ADC, Arduino UNO, Python, Numpy \& Scipy, Matplotlib, Scikit-learn, Tkinter}
\begin{itemize} \itemsep 1pt
    \item Engineered an \textbf{lead-1 analog front-end} using the AD620 instrumentation amplifier and filter stages (Right-leg drive, Notch, butterworth) to acquire low-amplitude (~1mV) heart signals amplified 1000x in ~V range using amplifier stages.
    \item Digitally sampled the analog out in 300-500 Hz, 16 bit precision using 16 bit ADC(ADS1115) with DC offset circuit (Clamper) and Arduino UNO as a data pipe to the processing computer (\textbf{client-side architecture)}.
    \item Digital signal undergoes segmentation, fragment selection extracting P-QRS-T feature space, then Principal Component Analysis (PCA) to compress 250-point cardiac windows into 30-point biometric "signatures". (\textbf{Feature space formation \& dimensionality reduction})
    \item Achieved a 96\% identification accuracy using Linear Discriminant Analysis (LDA) with Nearest Mean and majority vote classifiers on 30pt pca signature of 45 subjects (in ages 18-22) extracting 10 beats (P-QRS-T fragment) each, divided into 70-30 to train and predict, data acquired from public database. (\textbf{Classification \& Identification})
    \item Deployed a functional client-side testing environment featuring high-frequency Arduino sampling firmware and a custom \textbf{Tkinter-based interface} for real-time ECG recording and prediction. Architecting an edge-processing pipeline to migrate signal analysis from the client-side to Raspberry Pi/ESP32 platforms using \textbf{TinyML} for low-latency, on-device inference.
    \item Featured inherent \textbf{"Liveness Detection"} to prevent biometric spoofing common in fingerprint or face ID systems and hence proved uniqueness of normalised P-QRS-T ecg fragment across people, age and heart rate variability.
\end{itemize}
\vspace*{2mm}

\href{https://cerebralvalley.ai/e/google-deepmind-bangalore-hackathon/hackathon/gallery?project=81}{\textbf{Wake - On-Device AI Memory for Android}} \hfill \textit{Nov 2026}\\
{\sl Kotlin, Jetpack Compose (Material 3), Gradle, Room (SQLite/FTS4), MediaPipe TextEmbedder, LiteRT-LM, Gemini API}
\begin{itemize} \itemsep 1pt
    \item Built a local-first, on-device AI agent that captures useful text from Android notifications and the visible screen via an \textbf{Accessibility Service} and \textbf{Notification Listener Service}, storing it in a single \texttt{MemoryEvent} table -- no screenshots or pixels are ever captured -- and answers user queries about what was seen or received.
    \item Implemented an \textbf{agentic architecture}: \texttt{AgentReasoner} decides when to act, \texttt{CommitmentDetector} identifies user promises/intents from captured text, \texttt{ActionExecutor} executes actions, and \texttt{AgentEngine} orchestrates the perception $\to$ reasoning $\to$ action loop.
    \item Used \textbf{MediaPipe TextEmbedder} for on-device embeddings and \textbf{LiteRT-LM} for on-device \textbf{Gemma4:E2B} inference, with the Gemini API as an optional cloud LLM, in a single-module, no-DI, privacy-first architecture.
    \item Built at the \textbf{Google DeepMind $\times$ Cerebral Valley} hackathon (Bengaluru) using the latest Google AI stack (Gemini 3.5 Flash, Gemini Audio, GenMedia -- Omni, Nano Banana 2 Lite -- Interactions API/Antigravity, Gemma4); selected as one of the top 15-20 teams from 4,000+ applicants ($\sim$350 shortlisted).
\end{itemize}
\vspace*{2mm}

\textbf{Digital Logic Gate IC Tester} \hfill \textit{2026 -- Present}\\
{\sl AT89S52 (8051), Embedded C (SDCC/Keil C51), Arduino UNO, 16x2 LCD}
\begin{itemize} \itemsep 1pt
    \item Designing and building a low-cost microcontroller-based tester for verifying the functional operation of selected 74-series combinational logic ICs (\textbf{74LS00 NAND, 74LS02 NOR, 74LS04 NOT}), using an \textbf{AT89S52} 8051-family microcontroller as the central controller and a 14-pin ZIF socket to mount the Device Under Test (DUT).
    \item Push buttons provide user input for IC/test selection, while a \textbf{16x2 LCD} displays the selected IC and test status.
    \item Firmware applies the required input combinations, reads the corresponding outputs, and compares them against truth-table values stored in firmware to indicate pass/fail of the DUT.
    \item Validated the circuit through simulation before implementing on perfboard; firmware developed in \textbf{Embedded C} (SDCC or Keil C51) with an Arduino UNO used as the programming interface for the AT89S52.
\end{itemize}
\vspace*{2mm}

\href{https://github.com/aaromalonline/PassCrypt}{\textbf{PassCrypt - Password Manager}} \hfill \textit{2022}\\
{\sl Python, SQLite3}
\begin{itemize}
    \item[] A local password manager built with \textbf{Python and SQLite} for secure credential storage, using \textbf{SHA-256} hashing for user authentication and \textbf{AES encryption} (via cryptography.fernet) to keep stored passwords secure. Includes an advanced search/filter interface and a data-portability feature for exporting/importing the entire credential database.
\end{itemize}
\vspace*{2mm}

\href{https://github.com/aaromalonline/truevote}{\textbf{Truevote - Digital Voting System}} \hfill \textit{May 2025}\\
{\sl C++, Qt, SHA-256 Hashing, CMake}
\begin{itemize}
    \item[] A cross-platform digital voting application built with \textbf{C++ and Qt}, featuring an OOP-based admin/voter panel architecture for poll creation, voter registration, and authentication, with votes recorded and automatically tallied from CSV files. Uses \textbf{SHA-256} hashing for authentication and a \textbf{CMake}-driven build-and-run workflow for Linux.
\end{itemize}
\vspace{10pt}

%==== Experience ====%
\header{Experience}
\textbf{Technical Lead (R\&D AI/Robotics)} \hfill \textit{Jun 2026 -- Present} \\
\textit{\href{https://www.teamsingularit.com/}{Team SingularIT, CUSAT}} \hfill \textit{Kochi, India}
\begin{itemize}
    \item Leading R\&D on autonomous systems as Technical Lead for \textbf{Team SingularIT}, a student team preparing an autonomous entry for \textbf{A2RL} autonomous racing.
    \item Architected and developed the \textbf{perception subsystem} for the autonomous racing platform, running on an \textbf{NVIDIA Jetson} paired with an \textbf{STM32H743} ARM flight controller.
    \item Designed the perception stack to rely solely on \textbf{IMU and vision} (no external localization aids), fusing visual and inertial data for real-time state estimation.
\end{itemize}

\textbf{AI Research Intern -- Person Re-Identification} \hfill \textit{May 2026 -- Present} \\
\textit{\href{https://internship.iiitkottayam.ac.in/}{Indian Institute of Information Technology Kottayam (IIITK)}} \hfill \textit{Kottayam, India}
\begin{itemize}
    \item Conducting AI research under \textbf{Dr. Sridhar Raj S} on cloth-agnostic person re-identification (\textbf{CC-ReID, PU-ReID}), begun as a summer research internship under the Institute Internship Programme.
    \item Working on regularizing the attention of \textbf{Vision Transformers (ViT-Base, PyTorch)} to improve robustness to clothing changes.
    \item Wrote, trained, and evaluated vision models from scratch, running ablation studies that achieved improvements over current benchmark models. (\href{https://scholar.google.com/citations?user=dGwfp0wAAAAJ&hl=en}{Google Scholar})
\end{itemize}

\textbf{Software Intern} \href{https://github.com/aaromalonline/aaromalanil_horizon_s2}{\textit{(works)}}\hfill \textit{Mar 2025 -- Jul 2025} \\
\textit{\href{https://roverchallenge.eu/team/horizon-2024/}{Horizon Mars Rover Team, CUSAT}} \hfill \textit{Kochi, India}
\begin{itemize}
    \item Dual-booted \textbf{Ubuntu 22.04} by manual partitioning and setup \texttt{ros-humble-desktop-full} apt package.
    \item Architected node communication for a functional Mars rover using \textbf{ROS 2 Humble} and \textbf{Arduino framework}, wrote simple subscriber \& publisher nodes and proximity alert system using them in \textbf{Python \& C++}.
    \item Performed circuit simulation using \textbf{TinkerCAD} and \textbf{Arduino} to program robotic arm joints using servo-motors.
    \item Wrote technical reports of the rover for upcomming IROC and ERC competitions using \LaTeX\
    \item Track record for qualifying AIR 26 on IRC26' and global rank 18 on ERC24' rover challenges including the \textbf{European Rover Challenge (ERC)} and \textbf{IROC by ISRO}, India.
\end{itemize}

\textbf{Startup Founder | Founder Next Door (8-week mentorship program)} \hfill \textit{Oct 2025 -- Dec 2025} \\
\textit{\href{https://iedccusat.vercel.app/}{IEDC CUSAT, }}{\href{https://cittic.cusat.ac.in/}{CITTIC}} \hfill \textit{Kochi, India}
\begin{itemize}
    \item[] Selected from a campus-wide pool of proposals for an intensive 8-week incubator on market validation and MVP architecture. Designed the \textbf{system architecture} and business logic for \href{https://dozyo.vercel.app/}{\textbf{Dozyo}}, a hyperlocal platform for student-led gigs, and developed a comprehensive \href{https://aaromalonline.notion.site/Dozyo-1d91d429a96d80f18f9afc22c54f1637}{\textbf{business plan}} using Telegram community, Airtable, and escrow services to replace informal job-sharing with structured digital workflows. Collaborated weekly with industry founders on ideation, branding \& marketing, and technical scalability, concluding with an MVP presentation.
\end{itemize}
\vspace*{3mm}

%==== Volunteering ====%
\header{Volunteering}
\textbf{Active Member} \hfill \textit{Mar 2025 -- Present} \\
\textit{\href{https://fossunited.org/c/school-of-engineering-cusat}{FOSS CUSAT | Community}} \hfill \textit{Kochi, India}
\begin{itemize}
    \item[] Spearheaded the \textbf{\href{https://www.linkedin.com/posts/aaromalonline_endof10-foss-activity-7418182178624126976-I1s4?}{Linux Installation Party}} under the \textbf{\href{https://endof10.org/}{Endof10 Campaign}}, onboarding 20+ students to open-source operating systems with hands-on guidance on \textbf{manual disk partitioning}, Secure Boot configuration, and shell workflows across \textbf{Debian, Arch, Fedora, Ubuntu, and Pop!\_OS}, using deployment tools like \textbf{Ventoy} and \textbf{Balena Etcher}. Represented FOSS CUSAT at \textbf{\href{https://kochifoss.org/}{Kochi FOSS Meetups}} and \textbf{\href{https://fossunited.org/indiafoss/2025}{IndiaFOSS 2025}} (Bengaluru).
\end{itemize}
\vspace*{2mm}

\textbf{Community Volunteer} \hfill \textit{2026 -- Present} \\
\textit{\href{https://tinkerhub.org/tinkerspace}{Tinkerspace}, \href{https://tinkerhub.org/}{TinkerHub Foundation}} \hfill \textit{Kochi, India}
\begin{itemize}
    \item[] Host and actively engage in discussions at \textbf{Tinkerspace}, a local makerspace by TinkerHub, running \textbf{AI discussion sessions} and a \textbf{paper club}, mentoring at in-space hackathons like \textbf{Useless Projects}, and connecting with founders and fellow makers on their projects.
\end{itemize}
\vspace*{3mm}

\header{Technical Skills \& Certifications}
\begin{itemize} \itemsep 1pt
    \item \textbf{Cohere Labs Open Science Community ML Summer School '26} (Online):
    \begin{itemize}
        \item Attended sessions led by research scientists, PhD researchers, and mathematicians from \textbf{Meta Superintelligence Labs}, \textbf{Google DeepMind}, \textbf{Cohere Labs}, and \textbf{Liquid AI} on fundamental ML math, computer vision, and reinforcement learning.
        \item Covered vector calculus and linear algebra foundations of ML algorithms, the evolution of computer vision from classical techniques to transformer-based and vision-language architectures, and value-based/policy-based reinforcement learning.
    \end{itemize}
    \item \textbf{\href{https://matlabacademy.mathworks.com/?page=1&sort=featured}{MATLAB Academy} (Online Summer Internship):} \hfill \textit{Jun 2025}
    \begin{itemize}
        \item \href{https://www.credly.com/badges/5c9e05db-f4b5-4435-a360-226a04411a5e}{\textbf{MATLAB Proficiency:}} Core programming, algorithm development, and signal processing.
        \item \href{https://www.credly.com/badges/47d8690a-7a2d-4248-9c41-27fc8f8edc62}{\textbf{Data Analysis \& Processing:}} Visualization and AI/Machine Learning modules.
        \item \href{https://www.credly.com/badges/f69c1b8e-5a25-423f-9eba-db02fd2b6fb2}{\textbf{MATLAB for Simulink Users:}} System modeling, circuit simulation, and control logic.
    \end{itemize}

    \item \textbf{Languages:} Proficient in \textbf{Python} (NumPy, SciPy, Matplotlib, Seaborn, Scikit-learn, Pandas, OpenCV, Requests, Selenium, SQLite), \textbf{C/C++}, \textbf{x86 asm}, HTML/CSS, JS, Bash (basic scripting), SQL, \LaTeX\, and Verilog.
    \item \textbf{Tools \& Frameworks:} \textbf{ROS 2 Humble} (Basic Nodes), \textbf{Qt}, \textbf{Arduino} Framework, \textbf{LTspice}, \textbf{Xilinx Vivado}, Proteus, \textbf{PyTorch}, Kotlin/Jetpack Compose, and \textbf{MATLAB \& Simulink} (Basics), MySQL.
    \item \textbf{Linux \& DevOps:} \textbf{Git} workflow; Daily-driving \textbf{Linux} for $>$5 years across \textbf{Debian} (current), Ubuntu, Fedora, and Arch; experienced in VM management, manual SSD/HDD partitioning, Ventoy, and Balena Etcher.
\end{itemize}

\end{document}